% answers_short.tex —— 2005 年·普通高等学校招生考试·福建卷·文科 · 简短答案速查
% 由 tools/make_short_answers.py 自 answers.tex 的 \gksolhead/\gkanswer 提取，
% 并与 2005/2005_数学真题_福建文科_答案要点.md 逐题比对后生成，勿手改（改源 MD 后重跑）。

\begin{gknotice}
\textbf{使用说明}：本节只给最终答案便于快速核对，完整推导见解析版。答案为编写者独立推导，\textbf{非官方标准答案}。
\end{gknotice}

\gkpart{选择题}
\begin{center}\begin{tabular}{ccccccccccccc}
\toprule
题号 & 1 & 2 & 3 & 4 & 5 & 6 & 7 & 8 & 9 & 10 & 11 & 12 \\
\midrule
答案 & D & A & A & C & B & C & C & B & C & B & D & B \\
\bottomrule
\end{tabular}\end{center}

\gkpart{填空题}
\begin{itemize}[leftmargin=2.6em]
  \item[\textbf{13.}] 240
  \item[\textbf{14.}] 5
  \item[\textbf{15.}] 0
  \item[\textbf{16.}] 直线 $y = x$；$2^{x-3}$（或写关于 $x$ 轴对称，$g(x) = -3 - \log_2 x$ 等）
\end{itemize}

\gkpart{解答题}
\begin{itemize}[leftmargin=2.6em]
  \item[\textbf{17.}] (1) $\sin x - \cos x = -\dfrac{7}{5}$；(2) $\dfrac{\sin 2x + 2\sin^2 x}{1 - \tan x} = -\dfrac{24}{175}$
  \item[\textbf{18.}] (1) $\mathbf{E}\xi = \dfrac{9}{10}$；(2) 这四次投球中至少一次命中的概率为 $\dfrac{91}{100}$
  \item[\textbf{19.}] (1) $q = 1$ 或 $q = -\dfrac{1}{2}$；(2) 当 $q = 1$ 时，对所有 $n \geqslant 2$ 均有 $S_n > b_n$；当 $q = -\dfrac{1}{2}$ 时，当 $2 \leqslant n \leqslant 9$ 时 $S_n > b_n$，当 $n = 10$ 时 $S_{10} = b_{10}$，当 $n \geqslant 11$ 时 $S_n < b_n$
  \item[\textbf{20.}] (1) 函数解析式为 $f(x) = x^3 - 3x^2 - 3x + 2$；(2) 单调递增区间为 $(-\infty, 1 - \sqrt{2})$ 和 $(1 + \sqrt{2}, +\infty)$；单调递减区间为 $(1 - \sqrt{2}, 1 + \sqrt{2})$
  \item[\textbf{21.}] (1) 证明见解析（由 $CB \perp AE$ 且 $BF \perp AE$ 导出 $AE \perp \text{\text{平面} } BCE$）；(2) 二面角 $B - AC - E$ 的大小为 $\arccos\dfrac{\sqrt{3}}{3}$；(3) 点 $D$ 到平面 $ACE$ 的距离为 $\dfrac{2\sqrt{3}}{3}$
  \item[\textbf{22.}] (1) 椭圆 $C$ 的方程为 $\dfrac{x^2}{6} + \dfrac{y^2}{2} = 1$；(2) 存在，所求直线 $m$ 的方程为 $x = -2$，或 $x - \sqrt{3}y + 2 = 0$，或 $x + \sqrt{3}y + 2 = 0$
\end{itemize}
